\documentclass[11pt]{article}
\usepackage[margin=1in]{geometry}
\usepackage{amsmath,amssymb,amsthm,mathtools}
\usepackage{microtype}
\usepackage{enumitem}
\input{glyphtounicode}
\pdfgentounicode=1
\pdfglyphtounicode{mu}{03BC}
\pdfglyphtounicode{bardbl}{2016}
\pdfglyphtounicode{parenleftbig}{0028}
\pdfglyphtounicode{parenrightbig}{0029}
\pdfglyphtounicode{parenleftbigg}{0028}
\pdfglyphtounicode{parenrightbigg}{0029}
\pdfglyphtounicode{parenlefttp}{239B}
\pdfglyphtounicode{parenleftbt}{239D}
\pdfglyphtounicode{parenrighttp}{239E}
\pdfglyphtounicode{parenrightbt}{23A0}
\pdfglyphtounicode{circleplusdisplay}{2A01}
\pdfglyphtounicode{summationdisplay}{2211}
\pdfglyphtounicode{productdisplay}{220F}
\usepackage[colorlinks=true,linkcolor=blue,citecolor=blue,urlcolor=blue]{hyperref}
\usepackage{accsupp}
\NewCommandCopy{\OriginalHookrightarrow}{\hookrightarrow}
\NewCommandCopy{\OriginalLongmapsto}{\longmapsto}
\NewCommandCopy{\OriginalUpbracefill}{\upbracefill}
\renewcommand{\hookrightarrow}{\mathrel{
  \BeginAccSupp{method=hex,unicode,ActualText=21AA}
  \OriginalHookrightarrow\EndAccSupp{}}}
\renewcommand{\longmapsto}{\mathrel{
  \BeginAccSupp{method=hex,unicode,ActualText=27FC}
  \OriginalLongmapsto\EndAccSupp{}}}
\renewcommand{\upbracefill}{
  \BeginAccSupp{method=hex,unicode,ActualText=23DF}
  \OriginalUpbracefill\EndAccSupp{}}
\hypersetup{pdftitle={Quadratic stabilization of the canonical Foulkes--Howe map},pdfauthor={Rewritten for readability from the OpenAI manuscript of the same title}}
\newtheorem{theorem}{Theorem}
\newtheorem{corollary}[theorem]{Corollary}
\newtheorem{proposition}[theorem]{Proposition}
\newtheorem{lemma}[theorem]{Lemma}
\theoremstyle{remark}
\newtheorem{remark}[theorem]{Remark}
\DeclareMathOperator{\Sym}{Sym}
\DeclareMathOperator{\Av}{Av}
\newcommand{\C}{\mathbb C}
\newcommand{\Sa}{\mathfrak S_a}
\newcommand{\D}[2]{D_{#1,#2}}
\newcommand{\PP}{\mathcal P}
\newcommand{\unverified}[1]{\par\smallskip\noindent\fbox{\parbox{0.95\linewidth}{\textbf{Unverified step.} #1}}\par\smallskip}
\title{Quadratic stabilization of the canonical Foulkes--Howe map}
\author{Rewritten for readability from the OpenAI manuscript of the same title}
\date{September 25, 2026}
\begin{document}
\maketitle
\begin{abstract}
For every finite-dimensional complex vector space $V$, we prove that the
canonical Foulkes--Howe map
$\Sym^b(\Sym^a V)\longrightarrow\Sym^a(\Sym^b V)$ is surjective whenever
$a\ge2$ and $b\ge a(a-1)$. This gives a quadratic stabilization bound
independent of $\dim V$.
\end{abstract}

\section{The result}\label{sec:result}

\subsection{Setting}

Throughout, $V$ is a finite-dimensional complex vector space,
$\Sym V=\bigoplus_{j\ge0}\Sym^jV$ is its symmetric algebra, and a product
of vectors always means a product in $\Sym V$. Fix an integer $a\ge1$;
$\Sa$ is the symmetric group on $a$ letters.

\paragraph{Averaging.}
For a complex vector space $W$, the group $\Sa$ acts on $W^{\otimes a}$ by
permuting the tensor factors; $(W^{\otimes a})^{\Sa}$ denotes the invariant
tensors. The \emph{averaging operator} is the linear map
$\Av\colon W^{\otimes a}\to W^{\otimes a}$ given on simple tensors by
\begin{equation}\label{eq:average}
 \Av(w_1\otimes\cdots\otimes w_a)
 =\frac1{a!}\sum_{\pi\in\Sa}
 w_{\pi(1)}\otimes\cdots\otimes w_{\pi(a)}.
\end{equation}
We use three elementary facts.
\begin{enumerate}[label=(Av\arabic*),leftmargin=3.2em]
\item $\Av$ takes values in $(W^{\otimes a})^{\Sa}$ and fixes every
invariant tensor. Hence $\Av$ is a projection onto $(W^{\otimes a})^{\Sa}$,
and the elements $\Av(w_1\otimes\cdots\otimes w_a)$ span
$(W^{\otimes a})^{\Sa}$.
\item The right side of \eqref{eq:average} does not change when the $w_i$
are permuted, so $\Av$ induces a linear map
$\overline{\Av}\colon\Sym^aW\to(W^{\otimes a})^{\Sa}$,
$w_1\cdots w_a\mapsto\Av(w_1\otimes\cdots\otimes w_a)$.
\item $\overline{\Av}$ is an isomorphism; its inverse is the restriction to
invariant tensors of the quotient map $q\colon W^{\otimes a}\to\Sym^aW$.
Indeed, $q$ is unchanged by permuting factors, so
$q(\Av(w_1\otimes\cdots\otimes w_a))=w_1\cdots w_a$, i.e.\
$q\circ\overline{\Av}=\mathrm{id}$. Conversely,
$\overline{\Av}\circ q=\Av$ on all of $W^{\otimes a}$, because both sides
are linear and agree on simple tensors by the definition in (Av2); so for
an invariant $\tau$, $\overline{\Av}(q(\tau))=\Av(\tau)=\tau$ by (Av1).
\end{enumerate}

\paragraph{The algebra $R$.}
The tensor power $(\Sym V)^{\otimes a}$ is a commutative algebra under
factorwise multiplication,
$(s_1\otimes\cdots\otimes s_a)(s_1'\otimes\cdots\otimes s_a')
=s_1s_1'\otimes\cdots\otimes s_as_a'$, and $\Sa$ acts on it by algebra
automorphisms. For $j\ge0$ put
\[
 R_j=\bigl((\Sym^j V)^{\otimes a}\bigr)^{\Sa},\qquad
 R=\bigoplus_{j\ge0}R_j\subseteq(\Sym V)^{\otimes a}.
\]
Products of invariants are invariant, and $R_iR_j\subseteq R_{i+j}$, so $R$
is a graded commutative algebra. By (Av3), applied to $W=V$ and to
$W=\Sym^bV$,
\[
 R_1\cong\Sym^aV,\qquad R_b\cong\Sym^a(\Sym^bV).
\]

\paragraph{The map.}
Let $b\ge0$. Since $R$ is commutative, $(z_1,\dots,z_b)\mapsto z_1\cdots z_b$
is a symmetric $b$-linear map $R_1^{\,b}\to R_b$, hence induces a linear map
$\Sym^bR_1\to R_b$. Under the two identifications above this is the
\emph{canonical Foulkes--Howe map}
\begin{equation}\label{eq:map}
 \mu_{a,b,V}:\Sym^b(\Sym^a V)\longrightarrow\Sym^a(\Sym^b V).
\end{equation}
Let $A\subseteq R$ be the subalgebra generated by $R_1$ and $A_b=A\cap R_b$
its degree-$b$ part, i.e.\ the span of all products of $b$ elements of
$R_1$. By construction, $\operatorname{im}\mu_{a,b,V}$ corresponds to
$A_b$. Concretely (see Lemma~\ref{lem:image}): $\mu_{a,b,V}$ sends the
product in $\Sym^b(\Sym^aV)$ of the $b$ elements $u_1^a,\dots,u_b^a\in\Sym^aV$
to the $a$-th power, formed in $\Sym^a(\Sym^bV)$, of the single element
$u_1\cdots u_b\in\Sym^bV$.

\subsection{Main theorem}

\begin{theorem}\label{thm:main}
For every integer $a\ge2$, every integer $b\ge a(a-1)$, and every
finite-dimensional complex vector space $V$, the map
$\mu_{a,b,V}$ in \eqref{eq:map} is surjective. Equivalently, $A_b=R_b$.
\end{theorem}

For $a=1$ we have $R=\Sym V$ and $\mu_{1,b,V}$ is the identity in every
degree. The bound $a(a-1)$ does not involve $\dim V$. It is a sufficient
bound; it is not claimed to be sharp.

\subsection{Proof at a glance}\label{sec:glance}

\paragraph{Key idea.}
Surjectivity is dualized into a vanishing statement about multilinear
forms, and the vanishing is proved by the pattern
\emph{``differentiate into a known zero, then cancel the derivatives''}.
The cancellation is possible because a directional derivative
$\D{Y}{U}=\sum_kY_k\,\partial/\partial U_k$ is \emph{injective} on the
space of polynomials that are homogeneous of degree $p$ in the vector
variable $U$ and homogeneous of degree $q$ in the vector variable $Y$,
whenever $p>q$ (Lemma~\ref{lem:shift}).

\paragraph{Steps.}
\begin{enumerate}[leftmargin=2em]
\item \emph{Dualize (Lemma~\ref{lem:criterion}).} Call $u_1\cdots u_b\in
\Sym^bV$ a \emph{product of $b$ vectors}. Then $\mu_{a,b,V}$ is surjective if
and only if the only symmetric $a$-linear form $T$ on $\Sym^bV$ with
$T(P,\dots,P)=0$ for all products $P$ of $b$ vectors is $T=0$.
\item \emph{Vanishing (Proposition~\ref{prop:vanishing}).} This is proved
for $b\ge a(a-1)$ by induction on $a$. Write $r=a-1$ and $m=b-r$; the
hypothesis $b\ge a(a-1)$ is exactly $m\ge r^2$. The goal is
$T(v_1^b,\dots,v_r^b,t^b)=0$ for all vectors, which suffices because
$b$-th powers span $\Sym^bV$ (Lemma~\ref{lem:powers}).
 \begin{itemize}[leftmargin=1.5em]
 \item \emph{Step 1 (factoring over $\C$).} For a product $Q$ of $m$
 vectors, $(x^r+ct^r)Q$ is a product of $b$ vectors for every $c\in\C$.
 The coefficient of $c$ in the hypothesis gives
 $T(x^rQ,\dots,x^rQ,\,t^rQ)=0$. Here every slot contains the same $Q$.
 \item \emph{Step 2 (first differential sequence).} To replace the last
 slot $t^rQ$ by $t^b$, start from $H_x=T(x^rQ,\dots,x^rQ,\,t^b)$ with
 $Q=y_1\cdots y_m$ and apply $\D{y_1}{t},\dots,\D{y_m}{t}$. Each operator
 trades one factor $t$ of the last slot for one $y_i$, and the end result
 is a multiple of the Step~1 expression, hence $0$. Before the $i$-th
 operator the $t$-degree is $b-i+1$ and the $y_i$-degree is $r$, so each
 operator is injective and $H_x=0$.
 \item \emph{Step 3 (induction).} For fixed $x,t$, the form
 $(Q_1,\dots,Q_r)\mapsto T(x^rQ_1,\dots,x^rQ_r,t^b)$ on $\Sym^mV$ has $r$
 slots and, by Step~2, vanishing diagonal on products. The induction
 hypothesis (which needs $m\ge r(r-1)$) makes it zero for
 \emph{independent} $Q_1,\dots,Q_r$.
 \item \emph{Step 4 (second differential sequence).} Start from
 $f=T(v_1^b,\dots,v_r^b,t^b)$ and apply $\D{x}{v_i}$ $r$ times for each
 $i$. This trades $r$ factors $v_i$ of slot $i$ for $x^r$, and the end
 result is a multiple of $T(x^rv_1^m,\dots,x^rv_r^m,t^b)$, which is $0$ by
 Step~3. All $r^2$ operators raise the degree of the \emph{same} variable
 $x$, so before each of them the $x$-degree is at most $r^2-1$ (it is $0$
 before the first and $r^2-1$ before the last); the $v_i$-degree is at
 least $m+1$. Since $m+1>r^2-1$, each operator is injective and $f=0$.
 \end{itemize}
\end{enumerate}

\paragraph{Where the bound comes from.}
The hypothesis enters only through $m\ge r^2$ (that is,
$b\ge r(r+1)=a(a-1)$). What each step uses is this. Step~2 uses no lower
bound on $m$: its degree gaps $m-i+1$ ($1\le i\le m$) are automatically
positive. Step~3 needs $m\ge1$ and $m\ge r(r-1)$, in order to invoke the
proposition for the pair $(r,m)$. Step~4 needs $m+1>r^2-1$. All of these
follow from $m\ge r^2$; in Step~4 the inequality holds with room to spare
(the gap is at least $2$). The quadratic size arises in Step~4, where $r^2$
derivatives accumulate on the one variable $x$, together with the inductive
requirement of Step~3. The proof establishes sufficiency only; nothing here
asserts that $b\ge a(a-1)$ is necessary.

\paragraph{Dependencies.}
Lemma~\ref{lem:powers} (powers span) $\Rightarrow$ Lemma~\ref{lem:image}
(description of $A_b$) $\Rightarrow$ Lemma~\ref{lem:criterion} (criterion).
Lemma~\ref{lem:shift} (injectivity) $\Rightarrow$ Corollary~\ref{cor:peel}
(chains of operators). Lemmas~\ref{lem:powers} and~\ref{lem:deriv}
(derivative rule) and Corollary~\ref{cor:peel} $\Rightarrow$
Proposition~\ref{prop:vanishing}. Lemma~\ref{lem:criterion} and
Proposition~\ref{prop:vanishing} $\Rightarrow$ Theorem~\ref{thm:main}
$\Rightarrow$ Corollary~\ref{cor:comparison}.

\paragraph{Changes of notation relative to the source.}
The source writes $E$ both for $\Sym^bV$ and, inside the proof of the
injectivity lemma, for the adjoint operator $\D{U}{Y}$ of $\D{Y}{U}$. Here
$\Sym^bV$ is always written out, and that adjoint operator is called $D'$;
it occurs only in the proof of Lemma~\ref{lem:shift}. The operators applied
in Proposition~\ref{prop:vanishing} are always written $\D{Y}{U}$ with
explicit blocks. The source's $P_a=(\Sym V)^{\otimes a}$ is written out, so
$P$ always denotes a product of $b$ vectors. Statements left unnumbered in
the source are now Lemmas~\ref{lem:powers}, \ref{lem:image},
\ref{lem:criterion}, \ref{lem:deriv} and Corollary~\ref{cor:peel}. The
source's Lemma~2, Proposition~3, Corollary~4 and Remark~5 are
Lemma~\ref{lem:shift}, Proposition~\ref{prop:vanishing},
Corollary~\ref{cor:comparison} and Remark~\ref{rem:scope} here. The
literature discussion of the source's introduction is collected in
Section~\ref{sec:context}; no proof uses it.

\section{The canonical image and diagonal vanishing}\label{sec:map}

This section turns surjectivity of $\mu_{a,b,V}$ into a vanishing problem.
An element $P=u_1\cdots u_b\in\Sym^bV$ with $u_i\in V$ is called a
\emph{product of $b$ vectors}; the vectors may repeat or be zero. Such
products span $\Sym^bV$, because $\Sym^bV$ is the image of $V^{\otimes b}$.

\begin{lemma}[Polarization; pure powers span]\label{lem:powers}
Let $k\ge1$ and $u_1,\dots,u_k\in V$. In $\Sym V$,
\begin{equation}\label{eq:polar}
 k!\,u_1\cdots u_k
 =\sum_{J\subseteq\{1,\ldots,k\}}(-1)^{k-|J|}
   \left(\sum_{j\in J}u_j\right)^k .
\end{equation}
Consequently the pure powers $u^k$, $u\in V$, span $\Sym^kV$.
\end{lemma}

\noindent\emph{Used in:} Lemma~\ref{lem:image} (with $k=a$) and Step~5 of
Proposition~\ref{prop:vanishing} (with $k=b$).

\begin{proof}
\emph{Step 1: the identity for formal variables.} Work in
$\C[z_1,\dots,z_k]$. By the multinomial theorem,
$(\sum_{j\in J}z_j)^k=\sum_\beta\binom{k}{\beta}z^\beta$, the sum over
exponent vectors $\beta$ with $|\beta|=k$ whose support
$S=\{j:\beta_j>0\}$ lies in $J$. (For $J=\emptyset$ the left side is
$0^k=0$ because $k\ge1$, and the sum over $\beta$ is empty, so the formula
holds there too.) So the coefficient of $z^\beta$ on the right of
\eqref{eq:polar} is
\[
 \binom{k}{\beta}\sum_{J\supseteq S}(-1)^{k-|J|}
 =\binom{k}{\beta}\sum_{i=0}^{s}\binom{s}{i}(-1)^{s-i}
 =\binom{k}{\beta}(1-1)^{s},\qquad s=k-|S|,
\]
where we wrote $J=S\sqcup J'$ and grouped by $i=|J'|$. This is $0$ unless
$s=0$, i.e.\ $S=\{1,\dots,k\}$. With $|\beta|=k$ that forces
$\beta=(1,\dots,1)$, for which $\binom{k}{\beta}=k!$. So the right side
equals $k!\,z_1\cdots z_k$.

\emph{Step 2: substitution.} $\Sym V$ is a commutative $\C$-algebra, so
there is an algebra homomorphism $\C[z_1,\dots,z_k]\to\Sym V$ with
$z_j\mapsto u_j$. Applying it to Step~1 gives \eqref{eq:polar}.

\emph{Step 3: spanning.} $\Sym^kV$ is spanned by products $u_1\cdots u_k$.
By \eqref{eq:polar} and $k!\ne0$ in $\C$, each such product is a linear
combination of $k$-th powers.
\end{proof}

\begin{lemma}[The image]\label{lem:image}
For integers $a,b\ge1$,
\[
 A_b=\operatorname{span}_{\C}
 \{P^{\otimes a}: P\text{ is a product of }b\text{ vectors}\}
 \subseteq R_b .
\]
\end{lemma}

\noindent\emph{Used in:} Lemma~\ref{lem:criterion}.

\begin{proof}
(1) For $u\in V$, formula \eqref{eq:average} gives
$\overline{\Av}(u^a)=u^{\otimes a}=:z(u)\in R_1$, since all $a!$ summands
are equal.
(2) By Lemma~\ref{lem:powers} with $k=a$, the $u^a$ span $\Sym^aV$; since
$\overline{\Av}$ is an isomorphism (Av3), the $z(u)$ span $R_1$.
(3) $A_b$ is spanned by products $z_1\cdots z_b$ with $z_i\in R_1$, and the
product is multilinear, so $A_b$ is spanned by the products
$z(u_1)\cdots z(u_b)$.
(4) By the factorwise multiplication rule,
$z(u_1)\cdots z(u_b)=(u_1\cdots u_b)^{\otimes a}$.
\end{proof}

Under $R_b\cong\Sym^a(\Sym^bV)$ the tensor $P^{\otimes a}$ corresponds to the
$a$-th power of $P$ in $\Sym^a(\Sym^bV)$; this is the concrete formula for
$\mu_{a,b,V}$ quoted in Section~\ref{sec:result}.

\begin{lemma}[Diagonal criterion]\label{lem:criterion}
Let $a,b\ge1$. The map $\mu_{a,b,V}$ is surjective if and only if every
symmetric $a$-linear form $T$ on $\Sym^b V$ satisfying $T(P,\ldots,P)=0$ for
all products $P$ of $b$ vectors is zero.
\end{lemma}

\noindent\emph{Used in:} the proof of Theorem~\ref{thm:main} (only the
direction ``vanishing $\Rightarrow$ surjective'').

\begin{proof}
Recall $R_b=((\Sym^bV)^{\otimes a})^{\Sa}$, a subspace of
$(\Sym^bV)^{\otimes a}$.

\emph{Step 1: forms on $R_b$ are symmetric forms on $\Sym^bV$.} For
$\lambda\in R_b^*$ define
\[
 T_\lambda(w_1,\ldots,w_a)
 =\lambda\!\left(\Av(w_1\otimes\cdots\otimes w_a)\right),
 \qquad w_i\in\Sym^bV .
\]
This is $a$-linear, and symmetric by (Av2). We claim
$\lambda\mapsto T_\lambda$ is a bijection from $R_b^*$ onto the symmetric
$a$-linear forms on $\Sym^bV$.
\emph{Injective:} if $T_\lambda=0$ then $\lambda$ vanishes on all
$\Av(w_1\otimes\cdots\otimes w_a)$, which span $R_b$ by (Av1).
\emph{Surjective:} given a symmetric $a$-linear $T$, let
$\widetilde T\colon(\Sym^bV)^{\otimes a}\to\C$ be the linear map with
$\widetilde T(w_1\otimes\cdots\otimes w_a)=T(w_1,\dots,w_a)$ and put
$\lambda=\widetilde T|_{R_b}$. Then
\[
 T_\lambda(w_1,\dots,w_a)
 =\frac1{a!}\sum_{\pi\in\Sa}T(w_{\pi(1)},\dots,w_{\pi(a)})
 =T(w_1,\dots,w_a)
\]
by symmetry of $T$.

\emph{Step 2: the diagonal.} For a product $P$ of $b$ vectors,
$\Av(P^{\otimes a})=P^{\otimes a}$, so
$T_\lambda(P,\ldots,P)=\lambda(P^{\otimes a})$.

\emph{Step 3: conclusion.} By Lemma~\ref{lem:image} and Step~2, $\lambda$
vanishes on $A_b$ if and only if $T_\lambda(P,\dots,P)=0$ for all products
$P$. A subspace $A_b$ of the finite-dimensional space $R_b$ equals $R_b$ if
and only if no nonzero $\lambda\in R_b^*$ vanishes on it. By Step~1 this
says: $A_b=R_b$ if and only if the only symmetric $T$ with vanishing
diagonal on products is $T=0$. Finally $A_b=R_b$ is the surjectivity of
$\mu_{a,b,V}$.
\end{proof}

\begin{remark}[Reading of the criterion; not used later]\label{rem:reading}
A vector of $V$ is a linear form on $V^*$, so a product $P=u_1\cdots u_b$
is a degree-$b$ form on $V^*$ that splits into linear factors. The function
$w\mapsto T(w,\dots,w)$ is a homogeneous polynomial of degree $a$ on the
vector space $\Sym^bV$, and the criterion tests it on split forms. So here
$b$ is the degree of the forms and $a$ the degree of the test polynomial,
whereas in the algebra $A$ the integer $a$ counts tensor factors and $b$ is
the grading degree. Lemma~\ref{lem:criterion} is the duality between these
two descriptions; compare the product parametrization in
\cite[Section~2]{RSW}.
\end{remark}

\section{Differential tools}\label{sec:differential}

The proof of Proposition~\ref{prop:vanishing} recovers a polynomial from a
sequence of its directional derivatives. This section supplies the
injectivity statement that lets each derivative be cancelled
(Lemma~\ref{lem:shift}, Corollary~\ref{cor:peel}) and the rule for
computing the derivatives (Lemma~\ref{lem:deriv}).

\subsection{Vector variables}

Assume $n=\dim V>0$ and fix a basis $e_1,\dots,e_n$ of $V$. A \emph{vector
variable} $U$ is a block of $n$ independent indeterminates
$U_1,\ldots,U_n$. Let $\PP$ be a polynomial ring over $\C$ in the
coordinates of finitely many vector variables (and possibly further
indeterminates). A polynomial is \emph{homogeneous of degree $p$ in $U$} if
each of its monomials has total degree $p$ in $U_1,\dots,U_n$; the zero
polynomial is homogeneous of every degree. For two \emph{different} blocks
$U,Y$ define the operator on $\PP$
\[
 \D{Y}{U}=\sum_{k=1}^n Y_k\frac{\partial}{\partial U_k}.
\]
It lowers the $U$-degree by one, raises the $Y$-degree by one, and leaves
the degree in every other block unchanged.

\subsection{Injectivity}

\begin{lemma}\label{lem:shift}
Let $p,q$ be nonnegative integers. On polynomials homogeneous of degree $p$
in $U$ and degree $q$ in $Y$, the operator $\D{Y}{U}$ is injective if
$p>q$. The polynomials may depend on any finite collection of additional
variables independent of $U$ and $Y$.
\end{lemma}

\noindent\emph{Used in:} Corollary~\ref{cor:peel}, hence Steps~2 and~4 of
Proposition~\ref{prop:vanishing}.

The proof below is self-contained and is the only argument the paper relies
on for this lemma. The source remarks that the statement is the
positive-weight injectivity in $\mathfrak{sl}_2$ weight strings
\cite[Section~11.1, Claim~11.4 and equation~(11.5)]{FultonHarris} and that
Ikenmeyer \cite[Section~5, Claim~5.1]{Ikenmeyer} uses a two-coordinate
weight shift after decomposing tensors by untouched positions. These are
comparisons, not ingredients.

\begin{proof}
\emph{Plan.} With respect to a suitable inner product, the adjoint of
$D=\D{Y}{U}$ is $D'=\D{U}{Y}$. The commutator $D'D-DD'$ acts on
bihomogeneous polynomials as the scalar $p-q$. This gives the norm identity
\eqref{eq:norm}, which forbids $Dh=0$ with $h\ne0$ when $p>q$.

\emph{Step 1: the inner product.} Let $X=(X_1,\dots,X_N)$ list all
indeterminates of $\PP$, including the additional ones. Define a Hermitian
form (linear in the first argument, conjugate-linear in the second) by
\begin{equation}\label{eq:fischer}
 \Bigl\langle\sum_\alpha c_\alpha X^\alpha,\sum_\beta d_\beta X^\beta
 \Bigr\rangle=\sum_\alpha\alpha!\,c_\alpha\overline{d_\alpha},
 \qquad\alpha!:=\prod_k\alpha_k!.
\end{equation}
Distinct monomials are orthogonal, $\|X^\alpha\|^2=\alpha!$, and
$\langle g,h\rangle=\overline{\langle h,g\rangle}$. Since
$\langle h,h\rangle=\sum_\alpha\alpha!\,|c_\alpha|^2$, the form is positive
definite. This is the Fischer, or apolar, inner product
\cite[Section~2, equation~(8) and Proposition~3]{AldazRender}.

\emph{Step 2: differentiation is adjoint to multiplication.} For every
indeterminate $X_k$ and all $g,h\in\PP$,
\begin{equation}\label{eq:adjoint}
 \langle\partial_{X_k}g,h\rangle=\langle g,X_kh\rangle .
\end{equation}
Both sides are linear in $g$ and conjugate-linear in $h$, so it suffices to
take $g=X^\alpha$, $h=X^\beta$. The left side is
$\alpha_k\langle X^{\alpha-e_k},X^\beta\rangle$, which equals
$\alpha_k(\alpha-e_k)!=\alpha!$ if $\alpha=\beta+e_k$ and $0$ otherwise
(including when $\alpha_k=0$). The right side is
$\langle X^\alpha,X^{\beta+e_k}\rangle$, the same number.

We also need the mirrored form. Apply \eqref{eq:adjoint} with the roles of
$g$ and $h$ exchanged: $\langle\partial_{X_k}h,g\rangle=\langle
h,X_kg\rangle$. Taking complex conjugates of both sides and using
$\overline{\langle u,v\rangle}=\langle v,u\rangle$ gives
\begin{equation}\label{eq:adjoint2}
 \langle g,\partial_{X_k}h\rangle=\langle X_kg,h\rangle .
\end{equation}

\emph{Step 3: $D^*=D'$.} Put $D'=\D{U}{Y}=\sum_kU_k\,\partial_{Y_k}$.
Then
\[
 \langle Dg,h\rangle
 =\sum_k\langle Y_k\,\partial_{U_k}g,h\rangle
 =\sum_k\langle\partial_{U_k}g,\partial_{Y_k}h\rangle
 =\sum_k\langle g,U_k\,\partial_{Y_k}h\rangle
 =\langle g,D'h\rangle,
\]
where the second equality is \eqref{eq:adjoint2} (for the indeterminate
$Y_k$, with $\partial_{U_k}g$ in place of $g$) and the third is
\eqref{eq:adjoint} (for $U_k$). Exchanging $g$ and $h$ and conjugating
gives $\langle D'g,h\rangle=\langle g,Dh\rangle$.

\emph{Step 4: the commutator.} Let
$N_U=\sum_k U_k\partial_{U_k}$ and $N_Y=\sum_k Y_k\partial_{Y_k}$. As
operators, $\partial_{Y_l}Y_k=\delta_{kl}+Y_k\partial_{Y_l}$ and
$\partial_{U_k}U_l=\delta_{kl}+U_l\partial_{U_k}$, while $U$-operators
commute with $Y$-operators. Hence
\begin{align*}
 D'D&=\sum_{k,l}U_l\,\partial_{Y_l}Y_k\,\partial_{U_k}
    =N_U+\sum_{k,l}U_lY_k\,\partial_{Y_l}\partial_{U_k},\\
 DD'&=\sum_{k,l}Y_k\,\partial_{U_k}U_l\,\partial_{Y_l}
    =N_Y+\sum_{k,l}Y_kU_l\,\partial_{U_k}\partial_{Y_l},
\end{align*}
and the two double sums agree. Therefore
\begin{equation}\label{eq:comm}
 D'D-DD'=N_U-N_Y .
\end{equation}

\emph{Step 5: the norm identity.} Let $h$ be homogeneous of degree $p$ in
$U$ and $q$ in $Y$. On a monomial, $N_U$ multiplies by the total
$U$-degree and $N_Y$ by the total $Y$-degree (Euler), so $N_Uh=ph$ and
$N_Yh=qh$. By Step~3 and \eqref{eq:comm},
\begin{equation}\label{eq:norm}
 \|Dh\|^2-\|D'h\|^2
 =\langle h,D'Dh\rangle-\langle h,DD'h\rangle
 =\langle h,(N_U-N_Y)h\rangle
 =(p-q)\|h\|^2 .
\end{equation}

\emph{Step 6: conclusion.} If $Dh=0$ and $p>q$, then \eqref{eq:norm} reads
$-\|D'h\|^2=(p-q)\|h\|^2$. The left side is $\le0$ and the right side is
$\ge0$, so both vanish, and $p-q>0$ gives $\|h\|=0$, i.e.\ $h=0$.
\end{proof}

The lemma applies to the \emph{whole} space of polynomials of a given
bidegree, with all other variables as spectators. So it can be applied
repeatedly along a chain of operators, provided the degree inequality is
checked before each one:

\begin{corollary}[Cancelling a chain of derivatives]\label{cor:peel}
Let $F_0\in\PP$ and let $D_1,\dots,D_N$ be operators of the form
$D_j=\D{Y^{(j)}}{U^{(j)}}$, where $U^{(j)}\ne Y^{(j)}$ are blocks of $\PP$.
Put $F_j=D_jF_{j-1}$. Suppose that for each $j$ the polynomial $F_{j-1}$ is
homogeneous of some degree $p_j$ in $U^{(j)}$ and some degree $q_j$ in
$Y^{(j)}$, with $p_j>q_j$. If $F_N=0$, then $F_0=0$.
\end{corollary}

\begin{proof}
Descending induction on $j$: if $F_j=0$, then $D_jF_{j-1}=0$, and
Lemma~\ref{lem:shift} (with $p=p_j$, $q=q_j$) gives $F_{j-1}=0$.
\end{proof}

\subsection{Evaluating a multilinear form on vector variables}

We make precise what an expression such as $T(x^rQ,\dots,x^rQ,t^b)$ means
when $t$ and the factors of $Q$ are vector variables, and how $\D{Y}{U}$
acts on it.

\paragraph{Conventions.}
Elements of $\PP\otimes_\C\Sym V$ are polynomials in the indeterminates with
coefficients in $\Sym V$; this is a commutative ring. A vector variable $U$
is identified with the element $\sum_kU_ke_k$. An $a$-linear form $T$ on
$\Sym^bV$ extends uniquely to a $\PP$-multilinear map
$(\PP\otimes\Sym^bV)^a\to\PP$:
\[
 T(f_1s_1,\dots,f_as_a)=f_1\cdots f_a\,T(s_1,\dots,s_a),
 \qquad f_i\in\PP,\ s_i\in\Sym^bV .
\]
So $T$ of arguments involving vector variables is an element of $\PP$, an
ordinary polynomial with complex coefficients.

\emph{Evaluation.} Assigning complex values to all indeterminates is a ring
homomorphism $\mathrm{ev}\colon\PP\to\C$; applied to coefficients it gives
a ring homomorphism $\PP\otimes\Sym V\to\Sym V$, also written
$\mathrm{ev}$. It sends a vector variable $U=\sum_kU_ke_k$ to the vector
$\sum_k\mathrm{ev}(U_k)e_k$, and every vector of $V$ arises for a suitable
choice of values. It is compatible with $T$:
\[
 \mathrm{ev}\bigl(T(g_1,\dots,g_a)\bigr)
 =T(\mathrm{ev}\,g_1,\dots,\mathrm{ev}\,g_a).
\]
Both sides are additive in each $g_i$, and for $g_i=f_is_i$ both equal
$\mathrm{ev}(f_1)\cdots\mathrm{ev}(f_a)\,T(s_1,\dots,s_a)$, the right side
by multilinearity of $T$.

\emph{Zero principle.} A polynomial $F\in\C[X_1,\dots,X_N]$ that vanishes
at every point of $\C^N$ is zero. Induction on $N$: write
$F=\sum_jc_jX_N^j$ with $c_j\in\C[X_1,\dots,X_{N-1}]$. For fixed values of
$X_1,\dots,X_{N-1}$ the resulting polynomial in $X_N$ has infinitely many
roots, so all its coefficients vanish; thus each $c_j$ vanishes at every
point of $\C^{N-1}$ and is zero by induction. Combined with evaluation: an
expression $T(\cdots)\in\PP$ in vector variables is the zero polynomial as
soon as it vanishes for all choices of actual vectors.

\emph{Derivatives.} $\D{Y}{U}$ acts on $\PP\otimes\Sym V$ through the
coefficients: $\D{Y}{U}(fs)=\D{Y}{U}(f)\,s$. An element is called
\emph{$U$-free} if its coefficients do not involve $U_1,\dots,U_n$.

\begin{lemma}[Derivative rule]\label{lem:deriv}
Let $U\ne Y$ be blocks of $\PP$, $D=\D{Y}{U}$, and $T$ an $a$-linear form
on $\Sym^bV$. Let $1\le j\le b$, let $w\in\PP\otimes\Sym^{b-j}V$ and
$g_1,\dots,g_{a-1}\in\PP\otimes\Sym^bV$ all be $U$-free. Then
\begin{equation}\label{eq:rule}
 D\,T(g_1,\dots,g_{a-1},\,U^jw)
 =j\,T(g_1,\dots,g_{a-1},\,U^{j-1}Yw).
\end{equation}
The same identity holds, with the same proof, when $U^jw$ stands in any
other slot and the remaining $a-1$ arguments are $U$-free; symmetry of $T$
is not needed. The elements $w,g_i$ may involve $Y$.
\end{lemma}

\noindent\emph{Used in:} formulas \eqref{eq:first-composition} and
\eqref{eq:second-composition}.

\begin{proof}
(a) \emph{Leibniz rules.} For $g,h\in\PP\otimes\Sym V$ we have
$D(gh)=D(g)h+gD(h)$: both sides are biadditive, and for $g=fs$, $h=f's'$
the claim is $D(ff')ss'=(D(f)f'+fD(f'))ss'$, the Leibniz rule in $\PP$.
Similarly
$D\,T(g_1,\dots,g_a)=\sum_iT(g_1,\dots,Dg_i,\dots,g_a)$: both sides are
additive in each $g_i$, and for $g_i=f_is_i$ the left side is
$D(f_1\cdots f_a)\,T(s_1,\dots,s_a)$ with $T(s_1,\dots,s_a)\in\C$ a
constant, which the Leibniz rule in $\PP$ expands into the right side.

(b) $D(U)=\sum_kD(U_k)e_k=\sum_kY_ke_k=Y$, and $D(g)=0$ for $U$-free $g$,
because $D$ only differentiates in $U_1,\dots,U_n$.

(c) By (a) and (b), $D(U^j)=jU^{j-1}Y$ and
$D(U^jw)=jU^{j-1}Yw+U^jD(w)=jU^{j-1}Yw$. In the Leibniz rule for $T$ every
term in which $D$ falls on a $U$-free argument vanishes by (b), whichever
slot holds $U^jw$. This leaves \eqref{eq:rule}.
\end{proof}

\section{Vanishing on products}\label{sec:vanishing}

\begin{proposition}\label{prop:vanishing}
Let $V$ be a finite-dimensional complex vector space, and let $a\ge1$ and
$b\ge1$ be integers satisfying $b\ge a(a-1)$.
If a symmetric $a$-linear form $T$ on $\Sym^b V$ satisfies
\begin{equation}\label{eq:diagonal}
 T(P,\ldots,P)=0
 \qquad\text{for every product $P$ of $b$ vectors},
\end{equation}
then $T=0$.
\end{proposition}

\noindent\emph{Used in:} the proof of Theorem~\ref{thm:main}, via
Lemma~\ref{lem:criterion}.

\begin{proof}
\emph{Plan.} Induction on $a$, for all admissible $b$ and all $V$ at once.
In the inductive step we show $T(v_1^b,\dots,v_{a-1}^b,t^b)=0$ for all
vectors, via Steps~1--4 of Section~\ref{sec:glance}; Step~5 concludes by
spanning.

\paragraph{Trivial case and base case.}
If $V=0$ then $\Sym^bV=0$ (as $b\ge1$) and $T=0$. So let $n=\dim V>0$ and
fix a basis as in Section~\ref{sec:differential}. For $a=1$, $T$ is a
linear functional vanishing on all products of $b$ vectors; these span
$\Sym^bV$ (start of Section~\ref{sec:map}), so $T=0$.

\paragraph{Inductive setup.}
Let $a\ge2$ and assume the proposition for $a-1$ (all admissible degrees,
all $V$). Put
\begin{equation}\label{eq:rm}
 r=a-1\ge1,\qquad m=b-r.
\end{equation}
From $b\ge a(a-1)=r(r+1)$ we get
\begin{equation}\label{eq:range}
 b\ge r(r+1),\qquad m\ge r(r+1)-r=r^2\ge1.
\end{equation}

\paragraph{Step 1: a split binary form.}
\emph{Claim.} For all $x,t\in V$ and every product $Q=y_1\cdots y_m$ of $m$
vectors,
\begin{equation}\label{eq:binary}
 T(\underbrace{x^rQ,\ldots,x^rQ}_{r\text{ slots}},\,t^rQ)=0.
\end{equation}

(1a) \emph{$X^r+cY^r$ splits.} Let $c\in\C$. If $c=0$ the polynomial is
$X^r$. If $c\ne0$, choose $\alpha\in\C$ with $\alpha^r=-c$ and let
$\zeta=e^{2\pi i/r}$. In $\C[X,Y]$,
\[
 X^r+cY^r=X^r-\alpha^rY^r=\prod_{j=0}^{r-1}(X-\alpha\zeta^jY),
\]
because $X^r-\alpha^rY^r$ is monic of degree $r$ in $X$ over the domain
$\C[Y]$ and has the $r$ distinct roots $X=\alpha\zeta^jY$.

(1b) \emph{Substitution.} For any $x,t\in V$ (no independence needed) there
is an algebra homomorphism $\C[X,Y]\to\Sym V$ with $X\mapsto x$,
$Y\mapsto t$, because $\Sym V$ is commutative. It maps the factor
$X-\alpha\zeta^jY$ to the \emph{vector} $x-\alpha\zeta^jt\in V$. Hence
$x^r+ct^r$ is a product of $r$ vectors, and
\[
 P_c:=(x^r+ct^r)Q=x^rQ+c\,t^rQ
\]
is a product of $r+m=b$ vectors, for every $c\in\C$.

(1c) \emph{Coefficient of $c$.} By \eqref{eq:diagonal},
$T(P_c,\dots,P_c)=0$ for all $c\in\C$. Expanding by multilinearity,
$T(P_c,\dots,P_c)$ is a polynomial in $c$ of degree at most $a$. Its
coefficient of $c^1$ is the sum of the $a$ terms with $t^rQ$ in exactly one
slot and $x^rQ$ in the others; by symmetry of $T$ these are equal, giving
$a\,T(x^rQ,\ldots,x^rQ,t^rQ)$. A polynomial in $c$ vanishing for all
$c\in\C$ has all coefficients zero, and $a\ne0$. This proves
\eqref{eq:binary}.

\paragraph{Step 2: freeing the last slot.}
\emph{Claim.} For all $x,t\in V$ and every product $Q$ of $m$ vectors,
$T(x^rQ,\ldots,x^rQ,\,t^b)=0$.

Fix $x\in V$. Let $\PP$ be the polynomial ring in the coordinates of
independent vector variables $t,y_1,\ldots,y_m$, put $Q=y_1\cdots y_m$, and,
for $0\le k\le m$,
\[
 G_k=\frac{b!}{(b-k)!}\,
 T(\underbrace{x^rQ,\ldots,x^rQ}_{r\text{ slots}},\,t^{b-k}y_1\cdots y_k)
 \in\PP,\qquad H_x:=G_0=T(x^rQ,\ldots,x^rQ,\,t^b).
\]

(2a) \emph{$G_k=\D{y_k}{t}G_{k-1}$ for $1\le k\le m$}, that is,
\begin{equation}\label{eq:first-composition}
 \D{y_k}{t}\cdots\D{y_1}{t}H_x
 =\frac{b!}{(b-k)!}\,
 T(x^rQ,\ldots,x^rQ,\,t^{b-k}y_1\cdots y_k).
\end{equation}
Apply Lemma~\ref{lem:deriv} with $U=t$, $Y=y_k$, $j=b-k+1$ and
$w=y_1\cdots y_{k-1}$. The first $r$ arguments $x^rQ$ and $w$ are $t$-free,
so only the last slot is differentiated, and the scalar is
$\frac{b!}{(b-k+1)!}(b-k+1)=\frac{b!}{(b-k)!}$.

(2b) \emph{$G_m=0$.} Since $b-m=r$, we have
$G_m=\frac{b!}{r!}T(x^rQ,\ldots,x^rQ,\,t^rQ)$. At every assignment of
actual vectors to $t,y_1,\dots,y_m$ this vanishes by \eqref{eq:binary}, so
$G_m=0$ in $\PP$ by the zero principle.

(2c) \emph{Degrees before the $i$-th operator.} Let $1\le i\le m$. Under
$\PP$-multilinearity the degrees of the $a$ arguments add. Expanding
$y_l=\sum_ky_{l,k}e_k$ shows that each of the $r$ arguments $x^rQ$ has
degree $0$ in $t$ and degree $1$ in $y_i$ (the factor $y_i$ occurs in $Q$
exactly once). The last argument $t^{b-i+1}y_1\cdots y_{i-1}$ has degree
$b-i+1$ in $t$ and degree $0$ in $y_i$ (it contains only
$y_1,\dots,y_{i-1}$). Hence $G_{i-1}$ is homogeneous
\[
 \text{of degree }p=r\cdot0+(b-i+1)=b-i+1\text{ in }t
 \quad\text{and of degree }q=r\cdot1+0=r\text{ in }y_i,
\]
and $p-q=m-i+1\ge1$. The earlier operators
$\D{y_1}{t},\dots,\D{y_{i-1}}{t}$ raised the degrees of
$y_1,\dots,y_{i-1}$, not of $y_i$; this is why $q=r$ and not $r+1$.

(2d) \emph{Conclusion.} By (2a)--(2c), Corollary~\ref{cor:peel} applies to
$F_j=G_j$, $D_j=\D{y_j}{t}$, $N=m$, and gives $H_x=G_0=0$ in $\PP$.
Evaluating at actual vectors proves the claim.

\paragraph{Step 3: induction on the remaining slots.}
\emph{Claim.} For all $x,t\in V$,
\begin{equation}\label{eq:inductive-vanishing}
 T(x^rQ_1,\ldots,x^rQ_r,\,t^b)=0
 \qquad\text{for all }Q_1,\ldots,Q_r\in\Sym^m V.
\end{equation}

Fix $x,t\in V$ and define
$S_{x,t}(Q_1,\ldots,Q_r):=T(x^rQ_1,\ldots,x^rQ_r,t^b)$. Multiplication by
$x^r$ is a linear map $\Sym^mV\to\Sym^bV$ and $T$ is symmetric and
multilinear, so $S_{x,t}$ is a symmetric $r$-linear form on $\Sym^mV$. By
Step~2, $S_{x,t}(Q,\dots,Q)=0$ for every product $Q$ of $m$ vectors. The
induction hypothesis applies to $S_{x,t}$ with $(a,b)$ replaced by $(r,m)$:
indeed $r\ge1$, $m\ge1$, and $m\ge r^2\ge r(r-1)$ by \eqref{eq:range}.
Hence $S_{x,t}=0$.

\paragraph{Step 4: removing the common power.}
\emph{Claim.} $T(v_1^b,\ldots,v_r^b,t^b)=0$ for all
$v_1,\dots,v_r,t\in V$.

In \eqref{eq:inductive-vanishing} the $Q_i$ are independent; the only
remaining constraint is the common factor $x^r$. Fix $t\in V$. Let $\PP$
now be the polynomial ring in the coordinates of independent vector
variables $v_1,\ldots,v_r$ and one further vector variable $x$ (in
Steps~1--3, $x$ was a vector; here it is a variable). Put
\[
 f=T(v_1^b,\ldots,v_r^b,t^b)\in\PP ,
\]
which has $x$-degree $0$. Apply $r^2$ operators in the following order:
$\D{x}{v_1}$ $r$ times, then $\D{x}{v_2}$ $r$ times, \dots, then
$\D{x}{v_r}$ $r$ times. For $1\le i\le r$ and $0\le s\le r$ let $f_{i,s}$
be the result after all operators for $v_1,\dots,v_{i-1}$ and $s$ operators
for $v_i$ (so $f_{1,0}=f$ and $f_{i,r}=f_{i+1,0}$ for $i<r$).

(4a) \emph{Formula.} Recall $b-r=m$. For all such $i,s$,
\begin{equation}\label{eq:fis}
 f_{i,s}=\left(\frac{b!}{m!}\right)^{i-1}\frac{b!}{(b-s)!}\;
 T\bigl(x^rv_1^m,\ldots,x^rv_{i-1}^m,\;x^sv_i^{b-s},\;
 v_{i+1}^b,\ldots,v_r^b,\;t^b\bigr).
\end{equation}
We prove \eqref{eq:fis} by induction along the order in which the operators
are applied, i.e.\ on $(i,s)$ in lexicographic order.
\begin{itemize}[leftmargin=1.5em]
\item \emph{Start, $(i,s)=(1,0)$.} The right side of \eqref{eq:fis} is
$1\cdot\frac{b!}{b!}\,T(v_1^b,\dots,v_r^b,t^b)=f$.
\item \emph{From $(i,s)$ to $(i,s+1)$, for $0\le s\le r-1$.} Here
$f_{i,s+1}=\D{x}{v_i}f_{i,s}$. Apply Lemma~\ref{lem:deriv} (in its
``any slot'' form, for slot $i$) with $U=v_i$, $Y=x$, $j=b-s\ge m+1\ge1$
and $w=x^s\in\PP\otimes\Sym^sV$. The cofactor $x^s$ and the other slots
$x^rv_l^m$ ($l<i$), $v_l^b$ ($l>i$), $t^b$ are $v_i$-free: they involve
only $x$, the blocks $v_l$ with $l\ne i$, and the constant $t$. So only
slot $i$ changes, from $x^sv_i^{b-s}$ to $(b-s)\,x^{s+1}v_i^{b-s-1}$, and
the scalar becomes
$\frac{b!}{(b-s)!}(b-s)=\frac{b!}{(b-s-1)!}$. This is \eqref{eq:fis} for
$(i,s+1)$. In particular the previously inserted factors $x$ are not
differentiated.
\item \emph{From $(i,r)$ to $(i+1,0)$, for $i<r$.} These are the same
polynomial, and the two instances of \eqref{eq:fis} agree: for $(i,r)$,
slot $i$ is $x^rv_i^{b-r}=x^rv_i^m$ with scalar
$(b!/m!)^{i-1}\cdot b!/(b-r)!=(b!/m!)^{i}$; for $(i+1,0)$, slot $i+1$ is
$x^0v_{i+1}^b$ with scalar $(b!/m!)^{i}\cdot b!/b!$.
\end{itemize}
At the end, $(i,s)=(r,r)$ gives
\begin{equation}\label{eq:second-composition}
 \D{x}{v_r}^{\,r}\cdots\D{x}{v_1}^{\,r}f=f_{r,r}
 =\left(\frac{b!}{m!}\right)^r
 T(x^rv_1^m,\ldots,x^rv_r^m,t^b)=0.
\end{equation}
The last equality holds at every assignment of actual vectors by
\eqref{eq:inductive-vanishing} with $Q_i=v_i^m$, hence in $\PP$ by the zero
principle.

(4b) \emph{Degrees before each operator.} The operator applied to $f_{i,s}$
with $0\le s\le r-1$ is $\D{x}{v_i}$. Reading off \eqref{eq:fis} (degrees
of the arguments add): $v_i$ occurs only in slot $i$, with degree $b-s$;
and $x$ occurs with degree $r$ in each of the $i-1$ completed slots and
degree $s$ in slot $i$. So $f_{i,s}$ is homogeneous of degree $p=b-s$ in
$v_i$ and of degree $q=(i-1)r+s$ in $x$. Since $s\le r-1$ and $i\le r$,
\[
 p\ge b-r+1=m+1,\qquad q\le(r-1)r+(r-1)=r^2-1,\qquad
 p-q\ge m-r^2+2\ge2>0,
\]
using $m\ge r^2$ from \eqref{eq:range}. (Only $p-q>0$ is needed.) For
example, for $r=2$ the four operators see
$(p,q)=(b,0),(b-1,1),(b,2),(b-1,3)$.

(4c) \emph{Conclusion.} By (4a), (4b) and \eqref{eq:second-composition},
Corollary~\ref{cor:peel} applies to the chain of $r^2$ operators and gives
$f=0$ in $\PP$. Evaluating at actual vectors, and recalling that $t\in V$
was arbitrary, proves the claim.

\paragraph{Step 5: conclusion.}
By Lemma~\ref{lem:powers} with $k=b$, the pure powers $v^b$ span
$\Sym^bV$. By Step~4 (recall $a=r+1$), $T$ vanishes whenever all $a$
arguments are pure powers; by multilinearity $T=0$. This completes the
induction.
\end{proof}

\paragraph{Sanity check: $a=2$.}
Here $r=1$, $m=b-1\ge1$, and the bound is $b\ge2$. Step~1 gives
$T(xQ,tQ)=0$. Step~2 applies $b-1$ operators to $T(xQ,t^b)$ and reaches
$b!\,T(xQ,tQ)=0$, with degree gaps $b-1,b-2,\dots,1$. Step~3 is the base
case: $Q_1\mapsto T(xQ_1,t^b)$ is linear and vanishes on products. Step~4
is the single operator $\D{x}{v_1}$ on $T(v_1^b,t^b)$, which gives
$b\,T(xv_1^{b-1},t^b)=0$; it is injective since $b>0$.

\section{Proof of the theorem and a corollary}\label{sec:comparison}

\begin{proof}[Proof of Theorem~\ref{thm:main}]
Let $a\ge2$ and $b\ge a(a-1)$; then $b\ge2$, so $a,b\ge1$. By
Proposition~\ref{prop:vanishing}, every symmetric $a$-linear form on
$\Sym^bV$ whose diagonal vanishes on products of $b$ vectors is zero. By
Lemma~\ref{lem:criterion}, $\mu_{a,b,V}$ is surjective.
\end{proof}

\begin{corollary}\label{cor:comparison}
For every finite-dimensional complex vector space $V$, every integer
$a\ge2$, and every integer $b\ge a(a-1)$, there exists a
$\mathrm{GL}(V)$-equivariant embedding
\[
 \Sym^a(\Sym^b V)\hookrightarrow\Sym^b(\Sym^a V).
\]
In particular, taking $a=6$ (so $a(a-1)=30$), there is such an embedding
$\Sym^6(\Sym^bV)\hookrightarrow\Sym^b(\Sym^6V)$ for every $b\ge30$; this is
what the source calls ``the sixth symmetric-power comparison for
$b\ge30$''.
\end{corollary}

\begin{proof}
(1) \emph{$\mu_{a,b,V}$ is equivariant.} $\mathrm{GL}(V)$ acts on $\Sym V$
by algebra automorphisms, $g\cdot(u_1\cdots u_j)=(gu_1)\cdots(gu_j)$, hence
on each $W=\Sym^jV$ and diagonally on $W^{\otimes a}$. The diagonal action
commutes with permuting tensor factors, so it preserves each $R_j$ and, by
\eqref{eq:average}, commutes with $\Av$. Therefore
$\overline{\Av}\colon\Sym^aW\to(W^{\otimes a})^{\Sa}$ is equivariant, and
so are the identifications $R_1\cong\Sym^aV$ and
$R_b\cong\Sym^a(\Sym^bV)$. The diagonal action also respects factorwise
multiplication in $(\Sym V)^{\otimes a}$, so the multiplication map
$\Sym^bR_1\to R_b$ is equivariant. Hence $\mu_{a,b,V}$ is equivariant; it
is surjective by Theorem~\ref{thm:main}.
(2) Finite-dimensional polynomial representations of $\mathrm{GL}(V)$ over
$\C$ are completely reducible
\cite[Chapter~I, Appendix~A, Sections~7--8]{Macdonald}. So
$\ker\mu_{a,b,V}$ has an invariant complement $M$ in $\Sym^b(\Sym^aV)$.
(3) The restriction $\mu_{a,b,V}|_M$ is an equivariant isomorphism onto
$\Sym^a(\Sym^b V)$. Its inverse followed by the inclusion
$M\subseteq\Sym^b(\Sym^a V)$ is the required embedding.
This shows existence only; the embedding is not identified with any
prescribed canonical map.
\end{proof}

\begin{remark}[Scope]\label{rem:scope}
This remark restates Theorem~\ref{thm:main}; it claims nothing beyond it
(in particular nothing about infinite-dimensional $V$).
Lemma~\ref{lem:shift} holds for any number $n$ of coordinates and any
spectator variables, and the degree comparisons in
Proposition~\ref{prop:vanishing} involve only $a$ and $b$. Together with
the finite-dimensional duality of Lemma~\ref{lem:criterion}, this is why the
same bound holds for every finite-dimensional $V$. The argument treats each
$V$ directly: it does not first prove the statement in one convenient
dimension and then transfer it (the source's ``stable-dimension
reduction''). The source names three features of $\C$ used by the written
proof: positivity of the Fischer Hermitian inner product
(Lemma~\ref{lem:shift}); splitting of binary forms, which rests on $r$-th
roots (Step~1); and invertibility of the integer factors $k!$, $a$,
$b!/(b-k)!$ and $(b!/m!)^r$ in the polarization, coefficient and
differentiation identities. This list is not claimed to be exhaustive; for
instance, the zero principle also uses that $\C$ is infinite.
\end{remark}

\section{Context from the literature (not used in the proofs)}
\label{sec:context}

Everything in this section is reported from the source with its citations.
None of it was checked for this rewrite, and the source does not reproduce
the cited statements; read each sentence as ``the source states that
\dots''. Nothing in Sections~\ref{sec:map}--\ref{sec:comparison} depends on
this section.

\paragraph{The map in the literature.}
In positive dimension, Raicu--Sam--Weyman \cite[Section~2]{RSW} formulate
the map as $\Sym^b(D^aV)\longrightarrow D^a(\Sym^bV)$, with divided powers
realized as invariant tensors. Transporting both sides by the averaging
isomorphism (Av3) gives \eqref{eq:map}.

\paragraph{Geometry.}
The affine Chow variety of decomposable $a$-forms consists of products of
$a$ linear forms in $\Sym^a(V^*)$. Its coordinate ring identifies with $A$,
whose normalization is $R$; here $b$ is the coordinate degree
\cite[Section~2, Theorem~2.3]{RSW}. See also Landsberg
\cite[Section~7.6, Propositions~7.11--7.12]{Landsberg}.

\paragraph{Stabilization.}
For fixed $a$ and $V$, a \emph{stabilization degree} is an integer
$b_0\ge0$ such that $\mu_{a,b,V}$ is surjective for every $b\ge b_0$. One
may ask for the least such $b_0$ in a fixed dimension, or for a bound valid
in all finite dimensions at once; Theorem~\ref{thm:main} is of the second
kind. Brion established eventual surjectivity \cite{Brion1993} and then
gave an effective bound involving both $a$ and the dimension
\cite[Theorem~3.3]{Brion}. Since $(R/A)_b=R_b/A_b$ is nonzero exactly when
$\mu_{a,b,V}$ is not surjective, eventual surjectivity says that $R/A$ has
only finitely many nonzero graded components, and the least stabilization
degree is one more than the highest nonzero degree of $R/A$, or zero if
$R=A$. Raicu, Sam, and Weyman bound the regularity of the Segre ring and
modules of covariants, including the normalization
\cite[Theorem~3.9]{RSW}; according to the source, their theorem does not
compute the top degree of the quotient $R/A$. The source also states that
Theorem~\ref{thm:main} answers positively the request for a polynomial
bound in \cite[Problem~7.19]{Landsberg}: in that reference's notation
$d=b$, $n=a$ and $w=\dim V$, and the bound reads $d\ge n(n-1)$, a
polynomial in $n$ not involving $w$, valid in particular when $w=n$. The
problem's wording is not reproduced in the source, so neither the
translation of notation nor the claim that the problem is answered can be
checked from the manuscript; both are reported, not verified. The present
argument does not determine the least stabilization degree.

\paragraph{Relation to the Foulkes conjecture.}
The ordinary Foulkes conjecture \cite{Foulkes1950} asks, for $a\le b$, for
a $\mathrm{GL}(V)$-equivariant injection
$\Sym^a(\Sym^b V)\hookrightarrow\Sym^b(\Sym^a V)$
\cite[Conjecture~2.10]{RSW}. Corollary~\ref{cor:comparison} provides one in
the range $b\ge a(a-1)$. The conjecture is about the existence of some
embedding, whereas Theorem~\ref{thm:main} is about one specified map. The
difference is visible at $a=b$: source and target of $\mu_{a,a,V}$ are the
same representation, so an embedding exists trivially, yet $\mu_{a,a,V}$
need not be invertible. Since source and target have equal dimension, a
nonzero kernel is the same as a failure of surjectivity. The source reports
that the canonical map has nonzero kernel at $(a,b)=(5,5)$ and $(6,6)$ in
suitable dimensions
\cite[arXiv version, Theorem~6(b),(c) and Remark~7]{CIM}; it does not say
which dimensions, and the claim was not checked here. Whatever the
dimensions are, there is no conflict with Theorem~\ref{thm:main}: the
theorem covers every $V$, but for $a=5$ only $b\ge20$ and for $a=6$ only
$b\ge30$.

\paragraph{Relation of the method to earlier work.}
Successive weight shifts also occur in Ikenmeyer's proof
\cite[Section~5]{Ikenmeyer} of McKay's propagation theorem \cite{McKay}.
That setting differs from ours in two ways. It concerns the canonical maps
in the \emph{opposite} direction
$\Sym^a(\Sym^b V)\longrightarrow\Sym^b(\Sym^a V)$, and it assumes
$\dim V\ge b$ for each $b$. The result is a propagation statement: for
fixed $a$ and $B\ge a$, injectivity at $(a,B)$ propagates to every $(a,b)$
with $b\ge B$ \cite[Theorem~1.3 and Section~3]{Ikenmeyer}. Here there is no
dimension hypothesis; the two differential sequences (Steps~2 and~4)
surround an induction on the number of slots (Step~3) and give a
surjectivity range directly.

\section{Verification checklist}\label{sec:checklist}

No step in Sections~\ref{sec:map}--\ref{sec:comparison} carries an
``unverified'' mark, because no step was found that could not be justified.
That is a report, not a proof: the items below are the places where an
error would do the most damage and should be checked first. The literature
claims of Section~\ref{sec:context} were not checked and are not used.

\begin{enumerate}[leftmargin=2em]
\item \textbf{Duality} (Lemma~\ref{lem:criterion}, Steps~1--3): that
$\lambda\mapsto T_\lambda$ is a bijection, and that $A_b$ is spanned by
the tensors $P^{\otimes a}$ (Lemma~\ref{lem:image}, which uses
Lemma~\ref{lem:powers} with $k=a$).
\item \textbf{Injectivity lemma} (Lemma~\ref{lem:shift}): the adjointness
\eqref{eq:adjoint} and its mirrored form \eqref{eq:adjoint2}, the sign in
the commutator \eqref{eq:comm}, and the strict inequality $p>q$ in Step~6.
Also that $U$ and $Y$ must be different blocks.
\item \textbf{Step 1} of Proposition~\ref{prop:vanishing}: that
$(x^r+ct^r)Q$ really is a product of $b$ vectors for \emph{every} $c$
((1a)--(1b), using roots in $\C$), and that the coefficient of $c$ is
$a\,T(x^rQ,\dots,x^rQ,t^rQ)$ (1c).
\item \textbf{Step 2 degree count} (2c): before the $i$-th operator the
$t$-degree is $b-i+1$ and the $y_i$-degree is $r$, since $y_i$ occurs once
in each of the $r$ slots $x^rQ$ and not yet in the last slot; so
$p-q=m-i+1\ge1$.
\item \textbf{Step 3}: the induction hypothesis is applied to an $r$-linear
form on $\Sym^mV$, which requires $m\ge1$ and $m\ge r(r-1)$; both follow
from \eqref{eq:range}.
\item \textbf{Step 4 formula and degree count} ((4a), (4b)): the induction
proving \eqref{eq:fis}, including the hand-over from $(i,r)$ to $(i+1,0)$;
then that all $r^2$ operators raise the degree of the same block $x$, so
$q=(i-1)r+s$ can be as large as $r^2-1$, while $p=b-s\ge m+1$; the
inequality $m+1>r^2-1$ follows from $m\ge r^2$.
\item \textbf{Derivative rule} (Lemma~\ref{lem:deriv}): it requires only
that the other slots and the cofactor $w$ be free of the differentiated
variable $U$; they may contain $Y$, and the slot holding $U^jw$ is
arbitrary. This is what justifies \eqref{eq:first-composition} and
\eqref{eq:second-composition}.
\end{enumerate}

\begin{thebibliography}{99}
\bibitem{AldazRender}
J.~M.~Aldaz and H.~Render,
\href{https://doi.org/10.1007/s13324-023-00844-4}{\emph{A Fischer type decomposition theorem from the apolar inner product}},
Analysis and Mathematical Physics \textbf{13} (2023), article~91.

\bibitem{Brion1993}
M.~Brion,
\href{https://doi.org/10.1007/BF03026558}{\emph{Stable properties of plethysm: on two conjectures of Foulkes}},
Manuscripta Mathematica \textbf{80} (1993), no.~4, 347--371.

\bibitem{Brion}
M.~Brion,
\href{https://smf.emath.fr/publications/sur-certains-modules-gradues-associes-aux-produits-symetriques}{\emph{Sur certains modules gradu\'es associ\'es aux produits sym\'etriques}},
in \emph{Alg\`ebre non commutative, groupes quantiques et invariants}
(Reims, 1995), S\'eminaires et Congr\`es \textbf{2},
Soci\'et\'e Math\'ematique de France, Paris, 1997, 157--183.

\bibitem{CIM}
M.-W.~Cheung, C.~Ikenmeyer, and S.~Mkrtchyan,
\href{https://doi.org/10.1016/j.jsc.2016.09.002}{\emph{Symmetrizing tableaux and the 5th case of the Foulkes conjecture}},
Journal of Symbolic Computation \textbf{80} (2017), part~3, 833--843.
\href{https://arxiv.org/abs/1509.03944v1}{arXiv:1509.03944v1}.

\bibitem{Foulkes1950}
H.~O.~Foulkes,
\href{https://doi.org/10.1112/jlms/s1-25.3.205}{\emph{Concomitants of the quintic and sextic up to degree four in the coefficients of the ground form}},
Journal of the London Mathematical Society \textbf{25} (1950), no.~3,
205--209.

\bibitem{FultonHarris}
W.~Fulton and J.~Harris,
\href{https://doi.org/10.1007/978-1-4612-0979-9}{\emph{Representation Theory: A First Course}},
Graduate Texts in Mathematics \textbf{129}, Springer-Verlag, New York, 1991.

\bibitem{Ikenmeyer}
C.~Ikenmeyer,
\emph{On McKay's propagation theorem for the Foulkes conjecture},
\href{https://arxiv.org/abs/1509.04957v1}{arXiv:1509.04957v1}, 2015.

\bibitem{Landsberg}
J.~M.~Landsberg,
\href{https://doi.org/10.1007/s11565-014-0202-7}{\emph{Geometric complexity theory: an introduction for geometers}},
Annali dell'Universit\`a di Ferrara \textbf{61} (2015), no.~1, 65--117.
\href{https://arxiv.org/abs/1305.7387}{arXiv:1305.7387}.
Locators refer to the published version.

\bibitem{Macdonald}
I.~G.~Macdonald,
\href{https://doi.org/10.1093/oso/9780198534891.001.0001}{\emph{Symmetric Functions and Hall Polynomials}},
second edition, Oxford Mathematical Monographs, Clarendon Press, Oxford, 1995.

\bibitem{McKay}
T.~McKay,
\href{https://doi.org/10.1016/j.jalgebra.2007.12.003}{\emph{On plethysm conjectures of Stanley and Foulkes}},
Journal of Algebra \textbf{319} (2008), no.~5, 2050--2071.

\bibitem{RSW}
C.~Raicu, S.~V.~Sam, and J.~Weyman,
\href{https://doi.org/10.1007/s10013-021-00527-2}{\emph{On some modules supported in the Chow variety}},
Vietnam Journal of Mathematics \textbf{50} (2022), no.~2, 501--521.
\href{https://arxiv.org/abs/2108.10910v1}{arXiv:2108.10910v1}.
Locators refer to this preprint version.
\end{thebibliography}

\end{document}
